\section{PENDAHULUAN}

\lettrine{P}{asar} keuangan modern sangat dipengaruhi oleh dinamika nilai tukar, harga komoditas, dan tingkat suku bunga sehingga sering terdistorsi oleh korelasi antar aset, terutama selama periode krisis ekonomi global yang tidak terduga. Model \textit{Mean-Variance} (MV) yang diperkenalkan oleh Markowitz menjadi fondasi utama \textit{Modern Portfolio Theory} (MPT) untuk mengoptimasi keputusan investasi dengan menyeimbangkan ekspektasi imbal hasil dan risiko volatilitas. Namun, ketidakmampuan metode klasik dalam menangani ruang pencarian yang sangat luas mendorong perlunya transformasi diskritisasi fungsi objektif portofolio menjadi kerangka \textit{Quadratic Unconstrained Binary Optimization} (QUBO), yang memungkinkan pemanfaatan komputasi kuantum dalam mengidentifikasi konfigurasi aset optimal. Inisialisasi parameter yang murni acak pada algoritma \textit{Variational Quantum Eigensolver} (VQE) sangat rentan terjebak fenomena \textit{barren plateaus}, yaitu kondisi ketika gradien fungsi biaya meluruh secara eksponensial mendekati nol. Penelitian ini bertujuan mengevaluasi kemampuan algoritma hibrida GT-VQE yang mengintegrasikan ekuilibrium Nash dari \textit{Exact Potential Game} (EPG) sebagai strategi inisialisasi awal (\textit{warm-start}) guna memitigasi anomali \textit{barren plateaus}, mengakselerasi konvergensi VQE, dan menghasilkan keputusan alokasi aset yang lebih optimal pada sistem portofolio berskala $N=2$ dan $N=4$.

Kerangka teoritis pertama yang mendasari penelitian ini adalah \textit{Teori Permainan} (\textit{Game Theory}) dan \textit{Exact Potential Game} (EPG). Teori permainan merupakan kerangka kerja matematis aksiomatik yang digunakan untuk menganalisis interaksi strategis di antara sekumpulan agen rasional melalui tiga komponen fundamental: himpunan pemain ($N$), ruang strategi ($S$), dan fungsi utilitas ($U$). Kerangka EPG hadir untuk mengidentifikasi ekuilibrium Nash secara efisien, yakni kelas permainan strategis khusus di mana perubahan utilitas marginal setiap pemain selaras dengan variasi sebuah fungsi potensial global $\Phi(\mathbf{x})$. Dalam sistem pasar finansial, EPG berfungsi sebagai algoritma pencarian ekuilibrium Nash yang andal melalui mekanisme dinamika respons-terbaik (\textit{Best-Response Dynamics}) secara iteratif, di mana konfigurasi ekuilibrium biner yang diperoleh selanjutnya dimanfaatkan sebagai tebakan awal untuk memandu proses optimasi kuantum pada tahap berikutnya.

Kerangka teoritis kedua adalah \textit{Model Ising} dan penurunannya menuju Hamiltonian sebagai representasi sistem kuantum. Model Ising dikonseptualisasikan oleh Wilhelm Lenz dan diselesaikan oleh Ernst Ising pada tahun 1925 sebagai instrumen teoretis untuk menyelidiki asal-usul mikroskopis fenomena feromagnetisme, dengan memodelkan setiap partikel sebagai variabel spin biner $s_i \in \{-1,+1\}$ yang saling berinteraksi dalam kisi reguler. Meskipun solusi satu dimensi Ising gagal mendemonstrasikan transisi fase, paradigma ini berubah total setelah Lars Onsager merumuskan solusi eksak dua dimensi pada tahun 1944, yang membuktikan bahwa interaksi jarak pendek yang sederhana mampu memunculkan perilaku kolektif yang sangat kaya. Relevansi model Ising kemudian berkembang jauh melampaui magnetisme: dalam kerangka ekonofisika, setiap aset investasi dapat dianalogikan sebagai partikel spin, di mana korelasi antar pergerakan harga berperan sebagai koefisien kopling antar spin dan ekspektasi imbal hasil berperan sebagai pengaruh medan lokal. Transformasi inilah yang menjadi jembatan antara permasalahan optimasi portofolio diskret dan sistem fisik kuantum, karena fungsi energi total model Ising dapat diangkat menjadi operator Hamiltonian kuantum berbasis operator Pauli-$Z$ yang nilai eigen minimumnya --- yakni keadaan dasar (\textit{ground state}) --- langsung berkorespondensi dengan konfigurasi portofolio yang paling optimal.

Kerangka teoritis ketiga adalah metode optimasi parameter pada sistem klasik, yaitu \textit{Simultaneous Perturbation Stochastic Approximation} (SPSA). Algoritma SPSA digunakan karena mampu mengestimasi seluruh komponen gradien hanya menggunakan dua evaluasi fungsi biaya pada setiap iterasi, terlepas dari banyaknya parameter yang dimiliki sirkuit. Pada setiap iterasi ke-$k$, SPSA membangkitkan vektor perturbasi acak $\Delta_k$ berdistribusi Rademacher untuk membentuk dua titik evaluasi simetris $\boldsymbol{\theta}_k^{(\pm)} = \boldsymbol{\theta}_k \pm c_k\Delta_k$, yang menghasilkan estimator gradien tak-bias $\hat{g}_{k,i} = (E_k^{(+)} - E_k^{(-)})/(2c_k\Delta_{k,i})$. Karakteristik efisiensi evaluasi fungsi biaya yang konstan terhadap dimensi ruang parameter menjadikan SPSA sangat sesuai untuk optimasi sirkuit variasional berskala besar pada perangkat kuantum NISQ.

Kerangka teoritis keempat adalah algoritma VQE itu sendiri beserta arsitektur \textit{Ansatz} yang digunakan. VQE merupakan algoritma komputasi kuantum hibrida yang menggabungkan kemampuan pemrosesan kuantum dan optimasi klasik berdasarkan prinsip variasi Rayleigh-Ritz, yang menyatakan bahwa nilai ekspektasi energi $E(\boldsymbol{\theta}) = \langle\psi(\boldsymbol{\theta})|\hat{\mathcal{H}}|\psi(\boldsymbol{\theta})\rangle$ selalu merupakan batas atas energi keadaan dasar. Dalam penelitian ini, keadaan kuantum berparameter direpresentasikan menggunakan \textit{Hardware-Efficient Ansatz} EfficientSU(2), yang dibangun dari kombinasi gerbang rotasi satu-qubit ($\hat{R}_y$ dan $\hat{R}_z$) dan lapisan keterbelitan berbasis gerbang CNOT secara berulang. Ansatz ini dipilih karena menghasilkan sirkuit yang lebih dangkal dan lebih sesuai untuk perangkat kuantum era NISQ dibandingkan Ansatz berbasis \textit{Unitary Coupled Cluster} (UCC) yang memerlukan jumlah gerbang yang jauh lebih besar.

Kerangka teoritis kelima adalah validasi \textit{barren plateaus} menggunakan Entropi Von Neumann. Fenomena \textit{barren plateaus} merupakan kondisi ketika lanskap energi algoritma kuantum variasional menjadi sangat datar seiring bertambahnya jumlah qubit $N$ dan kedalaman sirkuit $L$, yang secara formal dapat dituliskan melalui hubungan $\text{Var}[\partial_\theta E] \approx C/2^{N+L}$. Untuk mendeteksi kondisi tersebut, penelitian ini memanfaatkan Entropi Von Neumann $S(\rho) = -\text{Tr}(\rho \log_2 \rho) = -\sum_j \lambda_j \log_2(\lambda_j)$ sebagai indikator tingkat keterbelitan dan distribusi probabilitas keadaan kuantum selama optimasi berlangsung. Peningkatan nilai entropi yang mendekati batas maksimumnya mengindikasikan bahwa keadaan kuantum menjadi semakin terdistribusi secara merata di dalam ruang Hilbert, sehingga sensitivitas Hamiltonian terhadap perubahan parameter menjadi sangat kecil yang secara empiris tercermin sebagai penurunan varians gradien pada proses optimasi.
